No auxiliary polynomial is introduced and no sign is imposed on
$\Delta_k^\Gamma$. The exact coordinate identification with (ACM1)--(ACM5)
is $y=u$, with $S=k/2+iu$ and unchanged ascending coefficients.
Its arithmetic form is the original $m_{h,k}=w_h^{*k}$, while the
reference form is precisely (ACM2), including its full convolution
mass. The relation matrices are the same multiplication by the full
$\chi$ of (HT2). Thus apply (ACM11a)--(ACM16) to this pair and write
$\beta_{p,\Gamma}^\pm$ for the resulting simultaneous endpoints,
whose initial value is $F_0=\mathcal B_k^\Gamma$. In particular
$\beta_{p,\Gamma}^-\leq\mathcal B_k\leq\beta_{p,\Gamma}^+$.
This uses every spectral, trace-norm, overlap, variance, angle and
constructed signed-residual control on the original Gamma-to-arithmetic
path. Put $T_k=4q\log(D_hk)$ for this scalar expression only. Then
(HC5), (HC8) and these endpoints give the current intervals
\[
 \boxed{
 \max\{0,-T_k,\beta_{p,\Gamma}^--T_k\}\leq\mathcal R_k
 \leq\min\{\widehat{\mathcal U}_k,\beta_{p,\Gamma}^+\}-T_k.}
 \tag{HC9}
\]
\[
\begin{split}
 \max\{0,T_k,\beta_{p,\Gamma}^-\}-\mathcal B_k^\Gamma
 &\leq\Delta_k^\Gamma\\
 &\leq\min\{\widehat{\mathcal U}_k,\beta_{p,\Gamma}^+\}
                          -\mathcal B_k^\Gamma.
\end{split}
\tag{HC10}
\]
Every entry follows by subtracting the same displayed initial value
from a proved bound on $\mathcal B_k$. Removing the source-dependent
$\beta_{p,\Gamma}^\pm$ entries gives the earlier coarse intervals.
Those coarse intervals have strictly positive width for every $k\geq3$.
